\section{Algorithms}
\label{sec:algorithms}

Reinforcement learning (RL) typically involves an agent interacting with an environment through states, actions, and rewards. In \emph{Q-learning}, the agent learns to approximate an optimal action-value function $Q(s,a)$, which represents the expected discounted reward for taking action $a$ in state $s$. By comparison, \emph{bandit} algorithms (including multi-armed and contextual bandits) consider no or minimal state transitions, learning directly which action maximizes the expected payoff.

\subsection{Q-Learning}

An \emph{asynchronous} Q-learning agent updates its $Q$-values only for the action actually taken at each step. Let $s_t$ be the current state, $a_t$ the chosen action, and $r_t$ the immediate reward upon transitioning to $s_{t+1}$. The Q-update is:
\begin{equation}
Q(s_t,a_t)
\;\leftarrow\;
Q(s_t,a_t)
\;+\;
\alpha \Bigl[
  r_t
  \;+\;
  \gamma \,\max_{a'} Q(s_{t+1},a')
  \;-\;
  Q(s_t,a_t)
\Bigr],
\label{eq:qlearning_async}
\end{equation}
where $\alpha$ is the learning rate and $\gamma$ the discount factor. Only $Q(s_t,a_t)$ changes, reflecting the experience from action $a_t$. A \emph{synchronous} variant updates Q-values for all actions simultaneously using counterfactual rewards, accelerating convergence at the cost of additional computation (Appendix~\ref{sec:appendix_algorithms}).

Agents select actions using either \emph{Boltzmann} (softmax) exploration, which draws actions probabilistically in proportion to exponentiated Q-values, or \emph{$\varepsilon$-greedy} exploration, which selects the greedy action with probability $1-\varepsilon$ and randomises otherwise. Both exploration methods use a decaying schedule that transitions from full exploration to near-greedy exploitation.\footnote{$\varepsilon$ starts at 1.0 and decays over the first 90\% of episodes to a floor of 0.01; agents then switch to pure exploitation ($\varepsilon = 0$) for the final 10\%. For Boltzmann exploration, the temperature $\tau$ follows an analogous schedule (1.0 to 0.01). Both schedules admit linear and exponential variants (Experiment~1).} Full specifications appear in Appendix~\ref{sec:appendix_algorithms}.

\subsection{Contextual Bandits}

\emph{Contextual bandits} generalise the multi-armed bandit problem by providing a context vector $\mathbf{x}\in\mathbb{R}^d$ prior to choosing an action. The \emph{LinUCB} algorithm assumes a linear payoff model, so each action $a$ has parameters $\boldsymbol{\theta}_a$, and the reward is approximately $\boldsymbol{\theta}_a^\top \mathbf{x}$. For each action $a$, one maintains a matrix $\mathbf{A}_a$ and vector $\mathbf{b}_a$. The algorithm selects $a$ via:
\begin{equation}
a_t
\;=\;
\arg\max_{a}\Bigl[\hat{\boldsymbol{\theta}}_{a}^\top \mathbf{x}
\;+\;
c\sqrt{\mathbf{x}^\top \mathbf{A}_a^{-1}\mathbf{x}}\Bigr],
\quad
\text{where}
\quad
\hat{\boldsymbol{\theta}}_{a} \;=\; \mathbf{A}_a^{-1}\mathbf{b}_a,
\end{equation}
and $c>0$ controls exploration. After observing the reward, $\mathbf{A}_a$ and $\mathbf{b}_a$ are updated, refining the estimate of $\boldsymbol{\theta}_a$. As an alternative, \emph{Contextual Thompson Sampling} (CTS) uses posterior sampling from a Bayesian linear model to balance exploration and exploitation; the two methods are compared in Experiment~3 (see Appendix~\ref{sec:appendix_algorithms} for the CTS formulation).

In standard implementations, the precision matrices $\mathbf{A}_a$ accumulate all past observations with equal weight, progressively reducing uncertainty and concentrating the policy around the estimated optimum. Experiment~3 introduces an optional \emph{memory decay} parameter $\delta \in (0,1]$ that applies exponential discounting to historical observations, preventing early experiences from permanently anchoring the policy.\footnote{At each update, $\mathbf{A}_a \leftarrow \delta\,\mathbf{A}_a + \mathbf{x}\mathbf{x}^\top$ and $\mathbf{b}_a \leftarrow \delta\,\mathbf{b}_a + r\,\mathbf{x}$. Setting $\delta = 1$ recovers the standard algorithm; $\delta < 1$ produces an effective observation window of approximately $1/(1-\delta)$ rounds.} This mechanism tests whether the rigidity of full-memory bandit learners contributes to bid suppression, following \citet{Douglas2024}, who show that deterministic convergence is a key driver of naive algorithmic collusion. The discounted variant draws on \citet{Russac2019}, who analyse weighted linear bandits in non-stationary environments.

\subsection{Budget-Constrained Pacing (Experiment~4a)}

Experiment~4a replaces the unconstrained bidding of prior experiments with \emph{budget-constrained pacing agents}. Each agent~$i$ has a total budget $B^i$ over a horizon of $T$ rounds. At each round~$t$, the agent observes its private valuation $v_t^i$ and computes a bid using a Lagrangian dual variable $\mu_t^i$ that controls bid shading. Following the dual decomposition framework of \citet{Balseiro2019Learning}, the bid depends on the agent's objective:
\begin{align}
b_t^i &\;=\; \min\!\bigl(v_t^i / \mu_t^i,\; B^i - S_t^i\bigr) & &\text{(value-maximizer),} \label{eq:bid_vmax_algo}\\
b_t^i &\;=\; \min\!\bigl(v_t^i / (1 + \mu_t^i),\; B^i - S_t^i\bigr) & &\text{(utility-maximizer),} \label{eq:bid_umax_algo}
\end{align}
where $S_t^i$ is the cumulative spend at round~$t$. The hard budget cap ensures spending never exceeds the remaining budget.

Following \citet{Balseiro2019Learning}, the dual variable is updated via an exponential multiplicative rule:
\begin{equation}
\mu_{t+1}^i \;=\; \operatorname{clip}\!\bigl(\mu_t^i \cdot \exp\!\bigl(\alpha_p\,(c_t^i - B^i/T)\bigr),\; \mu_{\min},\; \mu_{\max}\bigr),
\label{eq:dual_update}
\end{equation}
where $\alpha_p = 1/\!\sqrt{T}$ is the dual step size and $c_t^i$ is the payment made by agent~$i$ in round~$t$.\footnote{The dual variable is clipped to $[\mu_{\min}, \mu_{\max}] = [10^{-4}, 100]$ for numerical stability.} When the per-round cost exceeds the target spend rate $B^i/T$, the dual variable rises, reducing bids; underspending reverses this adjustment. An alternative PI-based controller is described in Appendix~\ref{sec:appendix_algorithms} and used in Experiment~4b.

Both algorithms implement the same closed-loop pacing structure illustrated in Figure~\ref{fig:pacing_loop}. The multiplier governs bids, the auction determines payments, and the control law updates the multiplier in response. They differ only in how the error signal is defined and propagated.

\begin{figure}[H]
\centering
\begin{tikzpicture}[
  node distance=1.8cm and 2.5cm,
  block/.style={rectangle, draw, rounded corners, minimum width=2.6cm,
                minimum height=0.9cm, text centered, font=\small},
  arrow/.style={-{Stealth[length=2mm]}, thick},
]
  % Nodes
  \node[block] (mult)   {Dual $\mu_t$};
  \node[block, right=of mult] (bid)    {Bid $b_t = \min(v_t/\mu_t,\,\text{rem})$};
  \node[block, right=of bid]  (auction){Auction};
  \node[block, below=1.2cm of auction] (cost)  {Cost $c_t$};
  \node[block, below=1.2cm of mult]  (ctrl)  {Control law};

  % Forward path
  \draw[arrow] (mult)    -- (bid);
  \draw[arrow] (bid)     -- (auction);
  \draw[arrow] (auction) -- (cost);
  \draw[arrow] (cost)    -- (ctrl);
  \draw[arrow] (ctrl)    -- node[left,font=\scriptsize]{$\mu_{t+1}$} (mult);

  % Labels on control arrow
  \node[font=\scriptsize, align=center, right=0.15cm of ctrl]
    {Mult: exp.\ update};
\end{tikzpicture}
\caption{Pacing feedback loop shared by both algorithms. The dual variable scales bids; auction outcomes feed back through the control law to update $\mu$.}
\label{fig:pacing_loop}
\end{figure}

The budget per episode is $B^i = 0.5 \cdot \mathbb{E}[v_t^i] \cdot T$, where $\mathbb{E}[v_t^i] = \exp(\mu_i + \sigma^2/2)$ under the log-normal valuation model. Note that $\eta$ throughout this paper denotes the affiliation parameter (Experiments~2 and~3); the dual step size above uses $\alpha_p$ to avoid notation collision.

In summary, Experiments~1--2 employ Q-learning under varying valuation models; Experiments~3a--3b replace Q-learning with contextual bandits (LinUCB and Thompson Sampling, respectively); and Experiments~4a--4b shift to budget-constrained pacing agents (multiplicative dual pacing and PI control, respectively).
